\section{Background}
\label{sec:background}

{\bf \em Code obfuscation}. Code obfuscation transforms a program while making its behavior harder to analyze. Obfuscation techniques vary in what they change and what they preserve. A classical taxonomy divides them into four families~\cite{collberg1997taxonomy,collberg2002manufacturing}: \emph{layout} transforms rewrite identifiers and remove names, comments, and annotations; \emph{control} transforms restructure the flow of computation; \emph{data} transforms change how values are represented or computed; and \emph{preventive} transforms target analysis tools rather than the program's representation. The first three are typically semantics-preserving, whereas preventive techniques can make behavior dependent on the execution environment through mechanisms such as anti-debugging checks, anti-virtualization probes, and self-modifying code.
Obfuscation techniques also operate at different levels. Source- and bytecode-level transformations leave an artifact that can be analyzed statically, whereas packers and virtualization-based techniques may reconstruct the protected code only at runtime~\cite{roundy2013packers,schrittwieser2016protecting}. Even when a transformation is semantics-preserving, preservation generally refers to observable program behavior: it does not require identical timing, memory usage, or exception behavior.

Our study focuses on semantics-preserving, source-level transformations. This restriction gives us a well-defined behavioral ground truth: because the transformation does not change program semantics, the behavior of an obfuscated program can be determined from its clean counterpart. This makes behavioral questions exactly gradable and allows us to test whether a model remains invariant to changes in program representation.

{\bf \em Predicting program behaviors}. We evaluate reasoning about obfuscated programs through two complementary behavior-prediction tasks: \emph{output prediction} and \emph{input prediction}. Given a program and an input, output prediction asks what the program returns. Conversely, given a program and a return value, input prediction asks for an input that produces that value. Current code models remain imperfect at both tasks, even for unobfuscated programs~\cite{cruxeval}.
%
The two directions probe different aspects of program reasoning. Output prediction follows the program's execution from input to result, whereas input prediction requires reasoning in the reverse direction. Evaluating both directions therefore provides a broader test of behavioral understanding than output prediction alone. This paired formulation is also motivated by prior work showing that reasoning performance can depend strongly on the direction represented during training~\cite{dexbench,revthink,cft_bidirectional,paired_attrib}.

This distinction is important for our study. If a model is trained to reason about one direction, improvements in that direction need not transfer to the other. We therefore treat the two tasks as separate but related capabilities, allowing us to measure not only whether a model can reason about obfuscated code, but also whether adaptation to one reasoning direction affects the other.

{\bf \em Adapting a model}. We use low-rank adaptation (LoRA)~\cite{lora} to adapt a code model to individual obfuscation transformations. LoRA freezes the base model weights $W$ and learns a low-rank update
$\Delta W = BA$, where $B \in \mathbb{R}^{d \times r}$, $A \in \mathbb{R}^{r \times k}$, and $r \ll d$. The resulting update can be stored separately from the base model and applied to the target modules.

This modularity is central to our study. Each LoRA adapter can be trained as a specialist for a particular obfuscation transformation, after which multiple specialists can be combined without changing the underlying base model. We keep the base model, target modules, and LoRA rank fixed across methods, so the primary differences between methods arise from how training data and learned updates are used rather than from model capacity.